\chapter{A máquina inteira}
% versão completa: chapters/10_synthesis.tex

\pencilfig{10}{\pencilart{10}}{A máquina inteira: o barramento de dados, os freios dentro dele e, em cima, quatro canais de rádio de controle.}

Pegamos um tipo de neurônio por vez e perguntamos o que ele acrescenta à máquina. É hora de juntar as peças.

O resultado é uma máquina de três camadas. Embaixo, a enorme rede telefônica dos neurônios \nt{GLUT}: por ela passam os dados. Ao lado, os neurônios inibitórios \nt{GABA}: eles controlam o fluxo, decidem o que passa, quando e com que volume. E acima de tudo isso, algumas estações de rádio, cada uma com o seu botão. \nt{DOPA} diz que mudanças de ligação manter. \nt{SERO} mantém o nível geral e, segundo a hipótese, o horizonte da paciência. \nt{NORA} escolhe entre explorar e decidir. \nt{ACET}, entre escrever e ler. Cada estação de rádio é, na verdade, um pequeno feixe de canais, não um fio só, mas mesmo assim os canais são um punhado para oitenta e seis bilhões de neurônios.

\section*{Uma história}

Vamos acompanhar um evento pela máquina inteira. O animal sabe há muito tempo: aperta a alavanca A, ganha suco. Um dia o mundo muda: A não funciona mais, e quem dá suco é a alavanca B, que o animal quase não usava.

O suco depois de A não veio: a dopamina responde com um silêncio, e as ligações que escolheram A enfraquecem. Os silêncios se repetem, os erros ficam maiores que o normal: a noradrenalina, segundo a hipótese, avisa que o mundo mudou. O botão de contraste desce, a escolha fica mais branda, e o ruído empurra com mais frequência a experimentar B. A acetilcolina passa a rede para o modo de escrita. Uma tentativa com B traz um suco inesperado: a dopamina dá uma rajada, as ligações que escolheram B se fortalecem. Depois de algumas tentativas, a expectativa alcança o mundo novo, as surpresas acabam, e os botões voltam ao lugar.

Repare: nenhum botão sabe o que os outros fazem. Cada um responde aos seus próprios sinais, e a ordem das ações se forma sozinha, como num circuito assíncrono em que cada bloco dispara quando as suas entradas estão prontas. Que cada botão, isoladamente, se comporta assim está em grande parte medido. Que eles trabalham juntos de forma coordenada, por enquanto, é hipótese.

\section*{Em que ela difere do seu notebook}

O notebook tem um clock comum; o cérebro, não. Os elementos do notebook são confiáveis; a sinapse perde mais da metade dos sinais. A memória do notebook fica separada do processador; no cérebro, ela está nas mesmas ligações. O notebook aprende (se aprende) por retropropagação do erro; o cérebro, por regras locais e por um único ``melhor do que o esperado'' difundido para todos. E o principal: o notebook gasta dezenas de watts num processador só, e o cérebro, vinte watts em tudo.

\section*{O que não sabemos}

Não sabemos quanto o córtex usa o tempo exato dos cliques. Não sabemos como ele ajusta bilhões de ligações em circuitos de muitos andares, se o professor é um só para todos. Não entendemos por completo o que as estações de rádio transmitem. Não sabemos como partes distantes do cérebro coordenam o trabalho sem um clock comum. E, por enquanto, ninguém sabe construir uma máquina que faça o mesmo com vinte watts.

Daqui em diante, o livro sai desse núcleo: vai para as entradas e saídas da máquina, para a sua depuração, para o sono e para as máquinas feitas de neurônios vivos num chip.

\fullversion{10}
